% Shared preamble for the study guides and mock exams (built by ../build.py).
% \sgroot is the path from the document's folder to study-guides/ (empty for guides, ../ for exams).
\providecommand{\sgroot}{}
\usepackage[margin=0.8in,headheight=15pt]{geometry}
\usepackage{fontspec}
\setmainfont{texgyrepagella}[Extension=.otf,UprightFont=*-regular,BoldFont=*-bold,ItalicFont=*-italic,BoldItalicFont=*-bolditalic]
\setsansfont{texgyreheros}[Extension=.otf,UprightFont=*-regular,BoldFont=*-bold,ItalicFont=*-italic,BoldItalicFont=*-bolditalic]
\usepackage{amsmath,amssymb,booktabs,longtable,array,calc,enumitem,xcolor,graphicx}
\usepackage{tikz}
\usepackage{pgfplots}\pgfplotsset{compat=1.18}
\usetikzlibrary{calc,arrows.meta,positioning}
\usepackage[most]{tcolorbox}
\usepackage{fancyhdr,titlesec,needspace,etoolbox,environ}
\usepackage{hyperref}

\definecolor{forest}{HTML}{24594B}
\definecolor{pale}{HTML}{EFF5F1}
\definecolor{keyc}{HTML}{C98A0B}   % key idea: amber
\definecolor{vocabc}{HTML}{2F6FA8} % vocabulary: blue
\definecolor{tipc}{HTML}{3F8A3A}   % make it make sense: green
\definecolor{watchc}{HTML}{B8452F} % watch out: red
\definecolor{linkc}{HTML}{7B4C9A}  % connection: purple
\hypersetup{colorlinks=true,linkcolor=forest,urlcolor=vocabc}

\pagestyle{fancy}\fancyhf{}
\renewcommand{\headrulewidth}{0pt}
\fancyfoot[C]{\small\sffamily\thepage}
\setlength{\parindent}{0pt}\setlength{\parskip}{5pt}
\setlength{\emergencystretch}{2em}
\setlist[itemize]{leftmargin=1.3em,itemsep=2pt,topsep=2pt}
\setlist[enumerate]{leftmargin=1.6em,itemsep=2pt,topsep=2pt}
\renewcommand{\arraystretch}{1.2}
\providecommand{\tightlist}{\setlength{\itemsep}{0pt}\setlength{\parskip}{0pt}}
\providecommand{\pandocbounded}[1]{#1}
\def\LTcaptype{none}

\titleformat{\section}{\Large\sffamily\bfseries\color{forest}}{}{0pt}{}[{\color{forest}\titlerule[0.8pt]}]
\titleformat{\subsection}{\large\sffamily\bfseries\color{forest}}{}{0pt}{}
\titleformat{\subsubsection}{\normalsize\sffamily\bfseries}{}{0pt}{}
\titlespacing*{\section}{0pt}{14pt}{8pt}
\titlespacing*{\subsection}{0pt}{12pt}{4pt}
\setcounter{secnumdepth}{0}

% Block quotes (textbook excerpts) as shaded boxes.
\renewenvironment{quote}{\begin{tcolorbox}[breakable,enhanced,colback=pale,colframe=forest,boxrule=0pt,leftrule=2.5pt,arc=0pt,
  left=8pt,right=8pt,top=4pt,bottom=4pt,before skip=4pt,after skip=6pt]\small}{\end{tcolorbox}}

% ------------------------------------------------------------------ annotated textbook pages
% Landscape letter page: the scanned page on the left (full height), notes on the right.
\newlength{\portraitW}\setlength{\portraitW}{8.5in}
\newlength{\portraitH}\setlength{\portraitH}{11in}
\def\pgmargin{21.6}  % pt, top/left margin around the scan
\def\pgheight{568.8} % pt, printed height of the scan (8.5in - 2 margins)
\newcommand{\setpagesize}[2]{%
  \global\pdfpagewidth=#1\global\pdfpageheight=#2%
  \global\paperwidth=#1\global\paperheight=#2}

\newcommand{\circlednum}[2]{\tikz[baseline=(c.base)]\node[circle,fill=#1,text=white,font=\scriptsize\sffamily\bfseries,inner sep=1.2pt,minimum size=11pt](c){#2};}

% \textbookpage{image}{page width pt}{page height pt}{lesson}{book page}{intro}{highlights}{notes}
% Notes (\addnote) are typeset into boxes first, then laid out: each sits level with its highlight
% when there is room, is pushed down to avoid the note above, and is pulled back up from the page bottom.
\newcount\nnotes
\newbox\hdrbox
\newdimen\prevbottom \newdimen\limitdim \newdimen\tmpdim
\def\notegap{7pt}
\newcommand{\addnote}[5]{% color, number, label, anchor (scan pt, <0 = none), text
  \global\advance\nnotes by 1
  \expandafter\ifx\csname nb@\the\nnotes\endcsname\relax\expandafter\newbox\csname nb@\the\nnotes\endcsname\fi
  \global\expandafter\setbox\csname nb@\the\nnotes\endcsname=\hbox{\tikz\node[notebox=#1]{\notehead{#1}{#2}{#3}#5};}%
  \expandafter\xdef\csname na@\the\nnotes\endcsname{#4}%
  \expandafter\xdef\csname nc@\the\nnotes\endcsname{#1}}
\newcommand{\nb}[1]{\csname nb@#1\endcsname}
\newcommand{\layoutnotes}{%
  \global\prevbottom=\dimexpr\pgmargin pt+\ht\hdrbox+\dp\hdrbox\relax
  \ifnum\nnotes>0
    \foreach \k in {1,...,\the\nnotes}{%
      \edef\a{\csname na@\k\endcsname}%
      \tmpdim=\dimexpr\prevbottom+\notegap\relax
      \ifdim \a pt<0pt\else
        \pgfmathsetlength{\dimen0}{\pgmargin+\a*\s-7}%
        \ifdim\dimen0>\tmpdim \tmpdim=\dimen0 \fi
      \fi
      \expandafter\xdef\csname ny@\k\endcsname{\the\tmpdim}%
      \global\prevbottom=\dimexpr\tmpdim+\ht\nb{\k}+\dp\nb{\k}\relax}%
    \global\limitdim=\dimexpr612pt-32pt\relax
    \foreach \k in {\the\nnotes,...,1}{%
      \tmpdim=\csname ny@\k\endcsname
      \ifdim\dimexpr\tmpdim+\ht\nb{\k}+\dp\nb{\k}\relax>\limitdim
        \tmpdim=\dimexpr\limitdim-\ht\nb{\k}-\dp\nb{\k}\relax
        \expandafter\xdef\csname ny@\k\endcsname{\the\tmpdim}%
      \fi
      \global\limitdim=\dimexpr\tmpdim-\notegap\relax}%
    \ifdim\limitdim<\dimexpr\pgmargin pt+\ht\hdrbox+\dp\hdrbox-\notegap\relax
      \typeout{WARNING: notes overflow on textbook page \currentbookpage}\fi
  \fi}
\newcommand{\drawnotes}{%
  \ifnum\nnotes>0
    \foreach \k in {1,...,\the\nnotes}{%
      \edef\a{\csname na@\k\endcsname}\edef\c{\csname nc@\k\endcsname}\edef\y{\csname ny@\k\endcsname}%
      \node[anchor=north west,inner sep=0pt] at ($(O)+(\notex pt,-\y)$) {\usebox{\nb{\k}}};
      \ifdim \a pt<0pt\else
        \draw[\c,line width=0.6pt] ($(scan.north east)+(0,{-\a*\s pt})$) -- ++(6pt,0) -- ($(O)+(\notex pt-1pt,-\y-7pt)$);
      \fi}%
  \fi}
\newcommand{\textbookpage}[8]{%
  \clearpage\setpagesize{11in}{8.5in}\thispagestyle{empty}%
  \def\currentbookpage{#5}%
  \pgfmathsetmacro{\s}{\pgheight/#3}%
  \pgfmathsetmacro{\pgright}{\pgmargin+#2*\s}%
  \pgfmathsetmacro{\notex}{\pgright+20}%
  \pgfmathsetmacro{\notew}{792-\notex-\pgmargin}%
  \pgfmathsetmacro{\notetw}{\notew-10}%
  \setbox\hdrbox=\vbox{\hsize=\notew pt\parskip=0pt
    {\raggedright\sffamily\bfseries\color{forest}#4\par}\vspace{1pt}
    {\sffamily\footnotesize\color{black!55}Miller \& Levine Biology, textbook page #5\par}\vspace{-1pt}
    {\color{forest}\rule{\notew pt}{0.6pt}}\par\vspace{2pt}
    \ifx\relax#6\relax\else{\small\raggedright #6\par}\fi}%
  \global\nnotes=0 #8\layoutnotes
  \null
  \begin{tikzpicture}[remember picture,overlay]
    \coordinate (O) at (current page.north west);
    \node[anchor=north west,inner sep=0pt,draw=black!25] (scan) at ($(O)+(\pgmargin pt,-\pgmargin pt)$)
      {\includegraphics[height=\pgheight pt]{#1}};
    \node[anchor=north west,inner sep=0pt] at ($(O)+(\notex pt,-\pgmargin pt)$) {\usebox{\hdrbox}};
    \node[anchor=south east,font=\scriptsize\sffamily,text=black!50,inner sep=0pt] at ($(O)+(792pt-\pgmargin pt,-603pt)$)
      {Part 1 · the textbook, annotated \quad·\quad Miller \& Levine Biology (2019), p.\,#5 \quad·\quad page \thepage};
    #7
    \drawnotes
  \end{tikzpicture}%
  \clearpage\setpagesize{\portraitW}{\portraitH}}

% Highlight one text line on the scan: \hl{color}{x0}{y0}{x1}{y1} in scan points (y down).
\newcommand{\hl}[5]{%
  \fill[#1,opacity=0.22] ($(scan.north west)+({(#2-1.5)*\s pt},{-(#3-1)*\s pt})$) rectangle
                         ($(scan.north west)+({(#4+1.5)*\s pt},{-(#5+1)*\s pt})$);}
% Numbered badge just left of a highlight.
\newcommand{\badge}[4]{%
  \node[circle,fill=#1,text=white,font=\scriptsize\sffamily\bfseries,inner sep=1pt,minimum size=10pt]
    at ($(scan.north west)+({#3*\s pt-7pt},{-#4*\s pt})$) {#2};}
\newcommand{\notehead}[3]{{\sffamily\bfseries\footnotesize\circlednum{#1}{#2}\ \color{#1}\MakeUppercase{#3}}\\[1pt]}
\tikzset{notebox/.style={text width=\notetw pt,inner sep=4pt,fill=#1!7,draw=#1!55,line width=0.4pt,
  font=\small,align=left,rounded corners=1.5pt}}

% ------------------------------------------------------------------ figures in the guide text
% Caption used under both book figures and new diagrams.
\newcommand{\figcaption}[2]{\par\smallskip{\small\sffamily\raggedright #1\par}%
  \ifx\relax#2\relax\else{\scriptsize\sffamily\color{black!55}Source: #2\par}\fi}
% \bookfig{png}{width fraction}{caption}{source}: a figure cropped from one of the books.
\newcommand{\bookfig}[4]{\par\Needspace{8\baselineskip}\medskip
  \begin{minipage}{\linewidth}\centering
    \fcolorbox{black!15}{white}{\includegraphics[width=#2\linewidth]{#1}}
    \figcaption{#3}{#4}
  \end{minipage}\par\medskip}
% Several book figures side by side: \bookfigrow{\bookfigcol{w}{png}{height}{caption}{source}\hfill ...}
\newcommand{\bookfigcol}[5]{\begin{minipage}[t]{#1\linewidth}\centering
    \fcolorbox{black!15}{white}{\includegraphics[height=#3,width=0.95\linewidth,keepaspectratio]{#2}}
    \figcaption{#4}{#5}\end{minipage}}
\newcommand{\bookfigrow}[1]{\par\Needspace{10\baselineskip}\medskip\noindent#1\par\medskip}
% New diagrams drawn for the guide: \begin{guidefig}...\end{guidefig}{caption}
\newenvironment{guidefig}[1]{\par\Needspace{10\baselineskip}\medskip\begin{minipage}{\linewidth}\centering\def\guidefigcap{#1}}
  {\figcaption{\guidefigcap}{}\end{minipage}\par\medskip}
\tikzset{
  gbox/.style={draw=#1,fill=#1!8,rounded corners=3pt,align=center,font=\small\sffamily,inner sep=5pt},
  garr/.style={-{Stealth[length=2.4mm]},line width=1pt,draw=#1},
  glab/.style={font=\scriptsize\sffamily,fill=white,inner sep=1.5pt,align=center},
}

% ------------------------------------------------------------------ figures shared by guides and exams
% \drawfig{name}: a TikZ/pgfplots figure in figures/name.tex
\newcommand{\drawfig}[1]{\input{\sgroot figures/#1}}
% \bookcrop[width]{id}{book}{page}{x0 y0 x1 y1}{label}{caption}: a figure cropped from a book page.
% book: C = Campbell Biology, ML = Miller & Levine Biology; page = printed page number;
% crop box in points from the page's top-left corner. build.py makes build/figs/<id>.png from these arguments.
\newcommand{\bookname}[1]{\ifstrequal{#1}{C}{Campbell Biology}{Miller \& Levine Biology}}
\newcommand{\bookcrop}[7][0.7]{\bookfig{\sgroot build/figs/#2.png}{#1}{#7}{\bookname{#3}, #6, p.\,#4}}
